% !TeX root = ../教材题库_全部条目_转写工作稿.tex
% Source: XJ-B1-C05-PR-P186-Y0375
% 第5章 / 5.3 三角函数的图象与性质 / 教材 p.186 / PDF 第193页
% OCR未校对：公式与文字均可能有误；未完成数学TeX转写；禁止直接发给学生
% 原题及相邻原解截图（核对用）：../题目原图/XJ-B1-C05-PR-P186-Y0375-01.webp
% 本文件只在转写完成后作为题干使用；OCR 条目尚不是可靠题干。
\begin{bankitem}{p.186 练习 第3题}{第5章 / 5.3 三角函数的图象与性质 / 教材 p.186 / PDF 第193页}{XJ-B1-C05-PR-P186-Y0375}{OCR未校对：公式与文字均可能有误；未完成数学TeX转写；禁止直接发给学生}
\begin{ocrdraft}
[教材p.186 / PDF第193页]
\par
3．利用函数的单调性，比较下列各组数的大小：
\par
(Dtan
\par
（2）tan
\par
tan
\end{ocrdraft}
\end{bankitem}
